\section{Cardinalidades e regras}

Para cada relacionamento, respondam: uma ocorrência de A se associa a quantas de B?
E o inverso? Justifiquem com BPMN, requisitos e regras do processo. Se não souberem,
registem a hipótese a validar.

\begin{table}[!ht]
\centering
\scriptsize
\setlength{\tabcolsep}{3pt}
\renewcommand{\arraystretch}{1.65}
\begin{tabularx}{\linewidth}{|>{\raggedright\arraybackslash}p{3.5cm}
                              |>{\centering\arraybackslash}X
                              |>{\centering\arraybackslash}X
                              |>{\raggedright\arraybackslash}X|}
\hline
\textbf{Relacionamento} &
\textbf{A $\rightarrow$ B} &
\textbf{B $\rightarrow$ A} &
\textbf{Regra de negócio / fonte} \\
\hline
\textbf{Estudante — realiza — Solicitação} &
1:N (hipótese) &
N:1 (hipótese) &
BPMN: atividade do estudante (enviar solicitação) \\
\hline
 & & & \\ \hline
 & & & \\ \hline
 & & & \\ \hline
 & & & \\ \hline
 & & & \\ \hline
 & & & \\ \hline
 & & & \\ \hline
 & & & \\ \hline
 & & & \\ \hline
 & & & \\ \hline
\end{tabularx}
\end{table}

\begin{warnbox}[Observação]
Se faltar linha, continue registrando os relacionamentos em uma folha adicional.
\end{warnbox}
